% ============================================================
%  CAPÍTULO 5 – DISCUSIÓN
% ============================================================

Los resultados presentados en el capítulo anterior confirman que el sistema CAIQ cubre los cinco objetivos específicos planteados con evidencia cuantitativa. Esta sección interpreta críticamente esos resultados en relación con las decisiones de diseño adoptadas, los sitúa frente al estado del arte y reflexiona sobre las aplicaciones específicas del sistema y los márgenes de mejora identificados.

\textbf{Justificación del paradigma de extracción basado en taxonomía.} La elección de una taxonomía controlada con reglas léxicas frente a modelos de reconocimiento de entidades (NER) se apoyó en la ausencia de un corpus de currículums anotados suficientemente grande. Esta decisión se alinea con los hallazgos de Zhang et al. (2022), que demuestran que sistemas apoyados en vocabularios bien estructurados pueden competir con modelos NER de mayor complejidad cuando la densidad de anotación es baja. Los resultados obtenidos validan esta elección: un \textit{recall} de 0,995 sobre el subconjunto de referencia manual ($n$\,=\,87) indica que la taxonomía captura prácticamente todas las competencias reales del candidato. El precio es una precisión moderada (0,698), consecuencia directa de la lógica de detección amplia: la taxonomía prioriza no omitir competencias reales (falsos negativos: solo 6) frente a no introducir competencias espurias (falsos positivos: 482). En el contexto de una herramienta de orientación profesional, esta asimetría es coherente con el objetivo del sistema, puesto que la penalización de un falso negativo —dejar de recomendar formación en una competencia que el candidato sí posee— es mayor que la de un falso positivo —sugerirle una formación adicional poco ajustada—. Enfoques como ESCOXLM-R (Zhang et al., 2023), que preentrenan modelos sobre la taxonomía ESCO, ofrecerían una alternativa más sofisticada a la coincidencia léxica, pero requieren infraestructura de inferencia incompatible con el despliegue sin GPU que impone el entorno de Hugging Face Spaces de este trabajo.

\textbf{Interpretación del rendimiento del recomendador y comparación con el estado del arte.} El estudio de ablación sobre $n$\,=\,482 candidatos revela que las cuatro variantes de \textit{scoring} convergen en un \textit{recall} de entre 0,406 y 0,413, con una diferencia máxima de 0,009 en la función objetivo compuesta. Esta convergencia tiene una doble lectura. Como hallazgo positivo, indica que el sistema es robusto frente a la elección de ponderación entre cobertura léxica y similitud semántica, lo que simplifica su mantenimiento y reduce el riesgo operativo asociado a futuras recalibraciones. Como señal crítica, apunta a que el techo de rendimiento actual no depende de la mezcla de componentes sino de la representación subyacente de competencias y recursos. Un \textit{recall} absoluto de 0,41 es coherente con los valores reportados en sistemas de \textit{person-job fit} de complejidad comparable (Qin et al., 2018) y refleja la dificultad estructural de recuperar recursos formativos relevantes cuando el espacio de competencias es escaso y la señal semántica entre el perfil del candidato y el contenido de los cursos es débil. La variante \texttt{hybrid\_tuned}, con $\alpha_{\text{cov}}=0{,}65$ y $\alpha_{\text{sem}}=0{,}35$, mantiene tanto la señal de cobertura léxica directa como la componente semántica que permite capturar afinidades entre competencias relacionadas, en línea con el marco híbrido de Burke (2002). Frente a enfoques de filtrado colaborativo como los de Koren et al. (2009) o He et al. (2017), CAIQ prescinde de la necesidad de historial de interacciones, lo que elimina el problema del arranque en frío (\textit{cold-start}) que limita la aplicabilidad de esos modelos a usuarios nuevos. Frente a aproximaciones basadas en grandes modelos de lenguaje (Hou et al., 2024; Zhao et al., 2023), CAIQ sacrifica flexibilidad a cambio de interpretabilidad plena y coste de inferencia nulo en producción.

\textbf{Limitaciones del corpus de datos y sostenibilidad del \textit{pipeline}.} La cobertura de competencias en el catálogo de cursos (15,7\,\% frente al 67,6\,\% de los másteres) constituye la limitación estructural más directa sobre la calidad de las recomendaciones de formación a corto plazo. La causa es la ausencia de descripciones de contenido en los cursos: solo el título está disponible como fuente textual, frente al campo \texttt{study\_content} que sí poseen los másteres. Ampliar la cobertura mediante extracción de descripciones extendidas, cuando estén disponibles en las plataformas, reduciría este \textit{gap} sin necesidad de cambios en la arquitectura del sistema. En cuanto al \textit{pipeline} de \textit{scraping}, los 187 ciclos ejecutados durante 71 días demuestran un comportamiento estable, pero el período de observación es insuficiente para confirmar la robustez ante cambios estructurales de los portales consultados (\texttt{Indeed}, \texttt{LinkedIn}, \texttt{Glassdoor}), cuya estructura HTML puede modificarse sin previo aviso. Este riesgo operativo es inherente a cualquier sistema que dependa de fuentes no estructuradas y deberá gestionarse mediante monitorización continua del \textit{pipeline} en producción.

\textbf{Aplicaciones específicas y alcance del sistema.} CAIQ está concebido para el profesional en transición hacia perfiles de datos, pero los resultados obtenidos abren la puerta a aplicaciones específicas en tres contextos distintos. En primer lugar, servicios de orientación universitaria y escuelas de negocio, donde el sistema puede integrarse como herramienta de diagnóstico previo al asesoramiento personalizado, reduciendo el tiempo de análisis y objetivando el punto de partida del candidato. En segundo lugar, departamentos de recursos humanos interesados en mapear el \textit{skill gap} interno de sus equipos de datos y generar planes de formación alineados con los requisitos reales del mercado, tal como recomienda da Silva et al. (2023) para los sistemas de recomendación educativa. En tercer lugar, la propia comunidad investigadora del mercado laboral de datos, para la que el corpus de 37.829 ofertas limpias y enriquecidas semánticamente con una taxonomía de 71 competencias constituye una fuente de análisis cuantitativo de la demanda, coherente con la taxonomía de seis familias de roles propuesta por Rahhal y Zein (2024).

En conjunto, los resultados muestran un sistema funcional, validado y desplegado, cuyas limitaciones identificadas —precisión del parser, techo de \textit{recall} del recomendador y cobertura del catálogo de cursos— son abordables mediante líneas de mejora concretas que se desarrollan en el capítulo siguiente.
